\chapter*{Introduction to Deep Learning}

\section*{What Is Deep Learning?}

Deep Learning is a branch of machine learning that enables computers to learn complex patterns directly from data using artificial neural networks with multiple layers. Unlike traditional machine learning methods, which often rely on manually designed features, deep learning automatically learns useful representations from raw data during training. This ability has significantly improved the performance of intelligent systems across a wide range of applications.

The term \textit{deep} refers to the presence of multiple hidden layers within a neural network. Rather than solving a problem in a single step, each layer transforms its input into a more meaningful representation before passing it to the next layer. As information flows through the network, increasingly complex patterns are extracted, allowing the model to solve tasks that would be difficult using conventional algorithms.

One of the key strengths of deep learning is its ability to perform \textbf{representation learning}. Instead of requiring domain experts to manually identify the most informative features, the network automatically discovers the features that best describe the data. For example, in image recognition, the first layers typically detect simple visual patterns such as edges, while deeper layers combine these patterns to recognize objects. Similar hierarchical representations are learned in natural language processing and speech recognition.

The rapid progress of deep learning has been driven by several technological advances, including the availability of large datasets, powerful Graphics Processing Units (GPUs), and improved optimization algorithms. Together, these advances have enabled researchers to train neural networks containing millions or even billions of parameters, leading to remarkable improvements in accuracy and scalability.

Today, deep learning is a central component of many modern artificial intelligence systems. It is widely used in computer vision, speech recognition, natural language processing, recommender systems, autonomous driving, robotics, medical image analysis, and generative artificial intelligence. Modern AI also draws on symbolic methods, probabilistic methods, optimization, search, reinforcement learning, and hybrid systems. Although these applications employ different architectures, they all rely on the same fundamental principle: learning through artificial neural networks.

Before exploring these advanced architectures, it is essential to understand the basic computational model on which they are built. The fundamental building block of every deep learning model is the artificial neuron, which forms the basis of artificial neural networks.
